\documentclass[10pt,a4paper]{amsart}
\usepackage{amsmath,amssymb,amsthm,graphicx}
\usepackage{verbatim}
\newenvironment{reference}{\comment}{\endcomment}
%\newenvironment{reference}{}{}
\newenvironment{slogan}{\comment}{\endcomment}
\newenvironment{history}{\comment}{\endcomment}
% For commutative diagrams we use Xy-pic
\usepackage[all]{xy}

% We use 2cell for 2-commutative diagrams.
\xyoption{2cell}
% Theorem environments.
%
\theoremstyle{plain}
\newtheorem{theorem}[subsection]{Theorem}
\newtheorem{proposition}[subsection]{Proposition}
\newtheorem{lemma}[subsection]{Lemma}

\theoremstyle{definition}
\newtheorem{definition}[subsection]{Definition}
\newtheorem{example}[subsection]{Example}
\newtheorem{exercise}[subsection]{Exercise}
\newtheorem{situation}[subsection]{Situation}

\theoremstyle{remark}
\newtheorem{remark}[subsection]{Remark}
\newtheorem{remarks}[subsection]{Remarks}

\numberwithin{equation}{subsection}

% Macros
%
\def\lim{\mathop{\mathrm{lim}}\nolimits}
\def\colim{\mathop{\mathrm{colim}}\nolimits}
\def\Spec{\mathop{\mathrm{Spec}}}
\def\Hom{\mathop{\mathrm{Hom}}\nolimits}
\def\Ext{\mathop{\mathrm{Ext}}\nolimits}
\def\SheafHom{\mathop{\mathcal{H}\!\mathit{om}}\nolimits}
\def\SheafExt{\mathop{\mathcal{E}\!\mathit{xt}}\nolimits}
\def\Sch{\mathit{Sch}}
\def\Mor{\mathop{\mathrm{Mor}}\nolimits}
\def\Ob{\mathop{\mathrm{Ob}}\nolimits}
\def\Sh{\mathop{\mathit{Sh}}\nolimits}
\def\NL{\mathop{N\!L}\nolimits}
\def\CH{\mathop{\mathrm{CH}}\nolimits}
\def\proetale{{pro\text{-}\acute{e}tale}}
\def\etale{{\acute{e}tale}}
\def\QCoh{\mathit{QCoh}}
\def\Ker{\mathop{\mathrm{Ker}}}
\def\Im{\mathop{\mathrm{Im}}}
\def\Coker{\mathop{\mathrm{Coker}}}
\def\Coim{\mathop{\mathrm{Coim}}}

% Boxtimes
%
\DeclareMathSymbol{\boxtimes}{\mathbin}{AMSa}{"02}

%
% Macros for moduli stacks/spaces
%
\def\QCohstack{\mathcal{QC}\!\mathit{oh}}
\def\Cohstack{\mathcal{C}\!\mathit{oh}}
\def\Spacesstack{\mathcal{S}\!\mathit{paces}}
\def\Quotfunctor{\mathrm{Quot}}
\def\Hilbfunctor{\mathrm{Hilb}}
\def\Curvesstack{\mathcal{C}\!\mathit{urves}}
\def\Polarizedstack{\mathcal{P}\!\mathit{olarized}}
\def\Complexesstack{\mathcal{C}\!\mathit{omplexes}}
% \Pic is the operator that assigns to X its picard group, usage \Pic(X)
% \Picardstack_{X/B} denotes the Picard stack of X over B
% \Picardfunctor_{X/B} denotes the Picard functor of X over B
\def\Pic{\mathop{\mathrm{Pic}}\nolimits}
\def\Picardstack{\mathcal{P}\!\mathit{ic}}
\def\Picardfunctor{\mathrm{Pic}}
\def\Deformationcategory{\mathcal{D}\!\mathit{ef}}
\usepackage{cleveref}
\crefname{originaldefinition}{original definition}{original definitions}
\Crefname{originaldefinition}{Original definition}{Original definitions}
\makeatletter
\newlabel{definition-coherent}{{definition-coherent}{source}{}{sourcepointer-1}{}}
\newlabel{definition-coherent@cref}{{[originaldefinition][0][]definition-coherent}{[1][0][]source}{}{}{}}
\makeatother
\crefname{originallemma}{original lemma}{original lemmas}
\Crefname{originallemma}{Original lemma}{Original lemmas}
\makeatletter
\newlabel{lemma-extension-finite-type}{{lemma-extension-finite-type}{source}{}{sourcepointer-2}{}}
\newlabel{lemma-extension-finite-type@cref}{{[originallemma][0][]lemma-extension-finite-type}{[1][0][]source}{}{}{}}
\makeatother
\crefname{originalreference}{original reference}{original references}
\Crefname{originalreference}{Original reference}{Original references}
\makeatletter
\newlabel{homology-lemma-characterize-weak-serre-subcategory}{{homology-lemma-characterize-weak-serre-subcategory}{source}{}{sourcepointer-3}{}}
\newlabel{homology-lemma-characterize-weak-serre-subcategory@cref}{{[originalreference][0][]homology-lemma-characterize-weak-serre-subcategory}{[1][0][]source}{}{}{}}
\makeatother
\title{The coherent modules form an abelian category and are stable under kernels, cokernels and extensions}
\author{The Stacks Project Authors}
\date{stacks-project:modules.tex:local-snapshot-20261012; GNU FDL 1.2 or later}
\begin{document}
\maketitle
\noindent\textit{Editorial note.} All delivered mathematical text is the author's own, copied from the pinned Stacks Project chapter source sources/stacks-modules/pinned-original-source.tex; no statement or proof step was rewritten, reordered or removed.\par\medskip
\noindent\textit{Editorial note.} Preamble provenance: the Stacks Project keeps one single shared preamble for every chapter; this card takes it directly from the staged sources/stacks-modules/preamble.tex as three verbatim spans (lines 8-12, 14-18 and 154-217), so no derived file was written and nothing was added or redefined. The wrapper class and the base packages are supplied by the assembler.\par\medskip
\noindent\textit{Editorial note.} Bibliography provenance: neither the counted statement nor its proof contains a citation (this lane rescans the whole span before writing the card), so the card carries no bibliography and no entry was reconstructed.\par\medskip
\noindent\textit{Editorial note.} Selection provenance: this card was opened under the file-based control contract (CONTROL.md), with the mechanical pre-card gate run on its statement and proof interval before the card was written.\par\medskip
\noindent\textit{Allowed standard prerequisites (editorial).}\begin{itemize}
\item The two clauses of the definition of a coherent module: local generation by finitely many sections and local finite generation of the kernels of the induced maps. Invoked in the counted proof (source lines 1881-1892) as a pointer to the definition of this chapter. Supplies both the statement that a subsheaf inherits and the finiteness input for the kernel of a map into a coherent sheaf.
\item The finiteness of extensions of modules of finite type, and the criterion for a category to be a weak Serre subcategory. Invoked inside the counted proof (source lines 1948 and 1979-1980) as a pointer to a lemma of this chapter and to a result of another chapter. Gives finite typeness of the middle term of a short exact sequence and the last assertion.
\end{itemize}\medskip
\ifdefined\hypertarget\hypertarget{sourcepointer-1}{}\fi
\noindent\textit{Source reference.} Original definition definition-coherent: sources/stacks-modules/pinned-original-source.tex, source lines 1881--1881. Navigation pointer only: the counted text refers to the Stacks item definition-coherent, which belongs to another chapter of the project and is not part of this card. The pointer carries the target's own label as its number because this chapter's staged snapshot does not contain the other chapter's tag table; its content is not reproduced here.\par\medskip
\ifdefined\hypertarget\hypertarget{sourcepointer-2}{}\fi
\noindent\textit{Source reference.} Original lemma lemma-extension-finite-type: sources/stacks-modules/pinned-original-source.tex, source lines 1667--1667. Navigation pointer only: the counted text refers to the Stacks item lemma-extension-finite-type, which belongs to another chapter of the project and is not part of this card. The pointer carries the target's own label as its number because this chapter's staged snapshot does not contain the other chapter's tag table; its content is not reproduced here.\par\medskip
\ifdefined\hypertarget\hypertarget{sourcepointer-3}{}\fi
\noindent\textit{Source reference.} Original reference homology-lemma-characterize-weak-serre-subcategory: sources/stacks-modules/pinned-original-source.tex, source lines 1980--1980. Navigation pointer only: the counted text refers to the Stacks item homology-lemma-characterize-weak-serre-subcategory, which belongs to another chapter of the project and is not part of this card. The pointer carries the target's own label as its number because this chapter's staged snapshot does not contain the other chapter's tag table; its content is not reproduced here.\par\medskip
\begin{lemma}
\label{lemma-coherent-abelian}
Let $(X, \mathcal{O}_X)$ be a ringed space.
\begin{enumerate}
\item Any finite type subsheaf of a coherent sheaf is coherent.
\item Let $\varphi : \mathcal{F} \to \mathcal{G}$
be a morphism from a finite type sheaf $\mathcal{F}$
to a coherent sheaf $\mathcal{G}$. Then $\Ker(\varphi)$ is of finite type.
\item Let $\varphi : \mathcal{F} \to \mathcal{G}$ be a morphism
of coherent $\mathcal{O}_X$-modules. Then
$\Ker(\varphi)$ and
$\Coker(\varphi)$ are coherent.
\item Given a short exact sequence of $\mathcal{O}_X$-modules
$0 \to \mathcal{F}_1 \to \mathcal{F}_2 \to \mathcal{F}_3 \to 0$
if two out of three are coherent so is the third.
\item The category $\textit{Coh}(\mathcal{O}_X)$ is a weak Serre subcategory
of $\textit{Mod}(\mathcal{O}_X)$. In particular, the category of
coherent modules is abelian and the inclusion functor
$\textit{Coh}(\mathcal{O}_X) \to \textit{Mod}(\mathcal{O}_X)$
is exact.
\end{enumerate}
\end{lemma}

\begin{proof}
Condition (2) of Definition \ref{definition-coherent}
holds for any subsheaf of a coherent sheaf. Thus we get (1).

\medskip\noindent
Assume the hypotheses of (2).
Let us show that $\Ker(\varphi)$ is of finite type. Pick $x \in X$.
Choose an open neighbourhood $U$ of $x$ in $X$ such
that $\mathcal{F}|_U$ is generated by $s_1, \ldots, s_n$.
By Definition \ref{definition-coherent} the kernel $\mathcal{K}$
of the induced map
$\bigoplus_{i = 1}^n \mathcal{O}_U \to \mathcal{G}$,
$e_i \mapsto \varphi(s_i)$ is of finite type.
Hence $\Ker(\varphi)$ which is the image of the
composition
$\mathcal{K} \to \bigoplus_{i = 1}^n \mathcal{O}_U \to \mathcal{F}$
is of finite type.

\medskip\noindent
Assume the hypotheses of (3).
By (2) the kernel of $\varphi$ is of finite type and
hence by (1) it is coherent.

\medskip\noindent
With the same hypotheses let us show that $\Coker(\varphi)$ is coherent.
Since $\mathcal{G}$ is of finite type so is $\Coker(\varphi)$.
Let $U \subset X$ be open and let $\overline{s}_i \in \Coker(\varphi)(U)$,
$i = 1, \ldots, n$ be sections. We have to show that the kernel of the
associated morphism
$\overline{\Psi} : \bigoplus_{i = 1}^n \mathcal{O}_U \to \Coker(\varphi)$
is of finite type. There exists an open covering
of $U$ such that on each open all the sections $\overline{s}_i$
lift to sections $s_i$ of $\mathcal{G}$. Hence we may assume
this is the case over $U$. We may in addition assume there are
sections $t_j$, $j = 1, \ldots, m$ of $\Im(\varphi)$ over $U$
which generate $\Im(\varphi)$ over $U$.
Let $\Phi : \bigoplus_{j = 1}^m \mathcal{O}_U \to \Im(\varphi)$
be defined using $t_j$ and
$\Psi :
\bigoplus_{j = 1}^m \mathcal{O}_U \oplus
\bigoplus_{i = 1}^n \mathcal{O}_U \to \mathcal{G}$
using $t_j$ and $s_i$.
Consider the following commutative diagram
$$
\xymatrix{
0 \ar[r] &
\bigoplus_{j = 1}^m \mathcal{O}_U \ar[d]_\Phi \ar[r] &
\bigoplus_{j = 1}^m \mathcal{O}_U \oplus
\bigoplus_{i = 1}^n \mathcal{O}_U \ar[d]_\Psi \ar[r] &
\bigoplus_{i = 1}^n \mathcal{O}_U \ar[d]_{\overline{\Psi}} \ar[r] &
0 \\
0 \ar[r] &
\Im(\varphi) \ar[r] &
\mathcal{G} \ar[r] &
\Coker(\varphi) \ar[r] &
0
}
$$
By the snake lemma we get an exact sequence
$\Ker(\Psi) \to \Ker(\overline{\Psi}) \to 0$. Since $\Ker(\Psi)$
is a finite type module, we see that $\Ker(\overline{\Psi})$ has finite type.

\medskip\noindent
Proof of part (4).
Let $0 \to \mathcal{F}_1 \to \mathcal{F}_2 \to \mathcal{F}_3 \to 0$
be a short exact sequence of $\mathcal{O}_X$-modules. By part
(3) it suffices
to prove that if $\mathcal{F}_1$ and $\mathcal{F}_3$ are coherent
so is $\mathcal{F}_2$. By Lemma \ref{lemma-extension-finite-type} we
see that $\mathcal{F}_2$ has finite type. Let
$s_1, \ldots, s_n$ be finitely many local
sections of $\mathcal{F}_2$ defined over a common open $U$ of $X$.
We have to show that the module of relations $\mathcal{K}$
between them is of finite type.
Consider the following commutative diagram
$$
\xymatrix{
0 \ar[r] &
0 \ar[r] \ar[d] &
\bigoplus_{i = 1}^{n} \mathcal{O}_U \ar[r] \ar[d] &
\bigoplus_{i = 1}^{n} \mathcal{O}_U \ar[r] \ar[d] &
0 \\
0 \ar[r] &
\mathcal{F}_1 \ar[r] &
\mathcal{F}_2 \ar[r] &
\mathcal{F}_3 \ar[r] &
0
}
$$
with obvious notation. By the snake lemma
we get a short exact sequence
$0 \to \mathcal{K} \to \mathcal{K}_3 \to \mathcal{F}_1$
where $\mathcal{K}_3$ is the module of relations among
the images of the sections $s_i$ in $\mathcal{F}_3$.
Since $\mathcal{F}_1$ is coherent we see that
$\mathcal{K}$ is the kernel of a map from a finite type module
to a coherent module and hence finite type by (2).

\medskip\noindent
Proof of (5). This follows because (3) and (4) show that
Homology, Lemma \ref{homology-lemma-characterize-weak-serre-subcategory}
applies.
\end{proof}

\end{document}
